% Source: https://artofproblemsolving.com/community/c3h2070676_quick_combinatorial_sum
% Stable identity: aops-topic-2070676-q01
Both sides are polynomials in $\alpha$, so it suffices to prove the identity for $0<\alpha<1$ (two polynomials agreeing at infinitely many points are equal).

Fix $0<\alpha<1$ and perform independent trials, each a success with probability $p=1-\alpha$ and a failure with probability $\alpha$. Consider the event $\mathcal E$: \emph{there are at least $n+1$ successes among the first $n+k+1$ trials.}

\emph{First count.} $\mathcal E$ happens exactly when the $(n+1)$-st success occurs at trial $n+1+j$ for some $0\le j\le k$. That happens when trial $n+1+j$ is a success and the previous $n+j$ trials contain exactly $n$ successes and $j$ failures, with probability $\dbinom{n+j}{j}p^{n+1}\alpha^j$. These events are disjoint, so
\[
P(\mathcal E)=p^{n+1}\sum_{j=0}^{k}\dbinom{n+j}{j}\alpha^j .
\]
\emph{Second count.} Directly, with $i=n+1+j$ successes among $n+k+1$ trials,
\[
P(\mathcal E)=\sum_{j=0}^{k}\dbinom{n+k+1}{n+1+j}p^{n+1+j}\alpha^{k-j}.
\]
Equating and dividing by $p^{n+1}>0$ gives the identity for $0<\alpha<1$, hence for all $\alpha$.

\emph{$\alpha=1$.} On the right only $j=0$ survives (as $(1-\alpha)^j=0$ for $j\ge1$), giving $\sum_{j=0}^k\dbinom{n+j}{j}=\dbinom{n+k+1}{n+1}=\dbinom{n+k+1}{k}$, the hockey-stick identity.

\emph{Limit.} For $0\le\alpha<1$, $P(\mathcal E)\to1$ as $k\to\infty$ (the number of successes in $n+k+1$ trials eventually exceeds $n$ with probability tending to $1$; e.g.\ for $k\ge n$ the probability of at most $n$ successes is at most $(n+1)\dbinom{n+k+1}{n}\alpha^{k+1}\to0$). Hence $p^{n+1}\sum_{j\ge0}\dbinom{n+j}{j}\alpha^j=1$, i.e.
\[
\boxed{\sum_{j=0}^{k}\dbinom{n+j}{j}\alpha^j=\sum_{j=0}^{k}\dbinom{n+k+1}{n+1+j}(1-\alpha)^j\alpha^{k-j}},\qquad \sum_{j=0}^{\infty}\dbinom{n+j}{j}\alpha^j=\dfrac{1}{(1-\alpha)^{n+1}}.
\]
(The case $\alpha=0$ of the limit is trivial.)
